%% kapa: exam=1 round=34 year=2023 period=1 form=A subject=civil_law subject_ko=민법 seq=32 number=32
%% source: Q-Net 감정평가사 기출문제 (cst00309) artlSeq=5211789 file=raw/1008/5211789/2023년도 제34회 감정평가사 1차시험 1교시.pdf
%% conversion: kordoc 4.24.51 -> tools/gen1.py (auto); review=proofread-sonnet; flags=none
\begin{question}{32}
부동산 실권리자명의 등기에 관한 법률에 대한 설명으로 옳은 것은? (다툼이 있으면 판례에 따름)
\begin{choices}
  \item 명의신탁자에게 법률효과를 직접 귀속시킬 의도의 매매계약을 체결한 사정이 인정되더라도, 부동산매매계약서에 명의수탁자가 매수인으로 기재되어 있다면 계약명의신탁으로 보아야 한다.
  \item 부동산소유권 또는 그 공유지분은 명의신탁 대상이 되지만, 용익물권은 명의신탁의 대상이 될 수 없다.
  \item 탈법적 목적이 없는 종중재산의 명의신탁에 있어서 종중은 명의신탁재산에 대한 불법점유자 내지 불법등기명의자에 대하여 직접 그 인도 또는 등기말소를 청구할 수 있다.
  \item 탈법적 목적이 없더라도 사실혼 배우자간의 명의신탁은 무효이다.
  \item 계약당사자인 매수인이 명의수탁자라는 사정을 매도인이 알지 못하였더라도, 매매로 인한 물권변동은 무효이다.
\end{choices}
\end{question}
\answer{4}
